\documentclass{article}

\usepackage{PRIMEarxiv}

\usepackage[utf8]{inputenc}
\usepackage[T1]{fontenc}
\usepackage{hyperref}
\usepackage{url}
\usepackage{booktabs}
\usepackage{amsfonts}
\usepackage{amsmath}
\usepackage{nicefrac}
\usepackage{microtype}
\usepackage{fancyhdr}
\usepackage{graphicx}
\usepackage{tikz}
\usetikzlibrary{arrows.meta,positioning,shapes.geometric,fit,calc,decorations.pathreplacing}
\usepackage{enumitem}
\usepackage{algorithm}
\usepackage{algpseudocode}
\usepackage{pifont}
\usepackage{amssymb}
\usepackage{xspace}
\usepackage[text=PREPRINT,color=gray!25,scale=3]{draftwatermark}
\graphicspath{{media/}}

\newcommand{\cmark}{\ding{51}}
\newcommand{\xmark}{\ding{55}}
\newif\ifanonymous
% \anonymoustrue
\anonymousfalse
% shared/macros.tex — Identity-masked macros for both arXiv and NDSS versions
% Toggle anonymization: set \anonymoustrue before % shared/macros.tex — Identity-masked macros for both arXiv and NDSS versions
% Toggle anonymization: set \anonymoustrue before % shared/macros.tex — Identity-masked macros for both arXiv and NDSS versions
% Toggle anonymization: set \anonymoustrue before \input{macros} for NDSS,
%                       set \anonymousfalse before \input{macros} for arXiv

\ifanonymous
  \newcommand{\toolname}{{\sc BoundaryFlow}\xspace}
  \newcommand{\anonrepo}{\url{https://anonymous.4open.science/r/XXXX}}
  \newcommand{\authorblock}{%
    \author{Anonymous Submission \\ Paper ID: TBD}%
  }
\else
  \newcommand{\toolname}{{\sc kcov-dataflow}\xspace}
  \newcommand{\anonrepo}{\url{https://github.com/yskzalloc/kcov-dataflow}}
  \newcommand{\authorblock}{%
    \author{Yunseong Kim \\
      Ericsson Software Technology \\
      \texttt{yunseong.kim@est.tech}}%
    % \author{Anonymous Submission}%
  }
\fi

% Technical shorthands
\newcommand{\sancov}{\texttt{SanitizerCoverage}\xspace}
\newcommand{\tracepc}{\texttt{trace-pc}\xspace}
\newcommand{\traceargs}{\texttt{trace-args}\xspace}
\newcommand{\traceret}{\texttt{trace-ret}\xspace}
\newcommand{\kcov}{\texttt{KCOV}\xspace}
\newcommand{\kasan}{\texttt{KASAN}\xspace}
\newcommand{\syzkaller}{\texttt{syzkaller}\xspace}
\newcommand{\drgn}{\texttt{drgn}\xspace}
\newcommand{\virtme}{\texttt{virtme-ng}\xspace}

% Data-flow tuple notation
\newcommand{\dftuple}[4]{\ensuremath{\langle #1, #2, #3, #4 \rangle}}
\newcommand{\callerid}{\ensuremath{\mathit{PC}_{\mathrm{caller}}}}
\newcommand{\calleeid}{\ensuremath{\mathit{PC}_{\mathrm{callee}}}}

% Listing style
\usepackage{listings}
\lstset{
  basicstyle=\ttfamily\footnotesize,
  breaklines=true,
  frame=single,
  numbers=left,
  numberstyle=\tiny,
  captionpos=b,
}
 for NDSS,
%                       set \anonymousfalse before % shared/macros.tex — Identity-masked macros for both arXiv and NDSS versions
% Toggle anonymization: set \anonymoustrue before \input{macros} for NDSS,
%                       set \anonymousfalse before \input{macros} for arXiv

\ifanonymous
  \newcommand{\toolname}{{\sc BoundaryFlow}\xspace}
  \newcommand{\anonrepo}{\url{https://anonymous.4open.science/r/XXXX}}
  \newcommand{\authorblock}{%
    \author{Anonymous Submission \\ Paper ID: TBD}%
  }
\else
  \newcommand{\toolname}{{\sc kcov-dataflow}\xspace}
  \newcommand{\anonrepo}{\url{https://github.com/yskzalloc/kcov-dataflow}}
  \newcommand{\authorblock}{%
    \author{Yunseong Kim \\
      Ericsson Software Technology \\
      \texttt{yunseong.kim@est.tech}}%
    % \author{Anonymous Submission}%
  }
\fi

% Technical shorthands
\newcommand{\sancov}{\texttt{SanitizerCoverage}\xspace}
\newcommand{\tracepc}{\texttt{trace-pc}\xspace}
\newcommand{\traceargs}{\texttt{trace-args}\xspace}
\newcommand{\traceret}{\texttt{trace-ret}\xspace}
\newcommand{\kcov}{\texttt{KCOV}\xspace}
\newcommand{\kasan}{\texttt{KASAN}\xspace}
\newcommand{\syzkaller}{\texttt{syzkaller}\xspace}
\newcommand{\drgn}{\texttt{drgn}\xspace}
\newcommand{\virtme}{\texttt{virtme-ng}\xspace}

% Data-flow tuple notation
\newcommand{\dftuple}[4]{\ensuremath{\langle #1, #2, #3, #4 \rangle}}
\newcommand{\callerid}{\ensuremath{\mathit{PC}_{\mathrm{caller}}}}
\newcommand{\calleeid}{\ensuremath{\mathit{PC}_{\mathrm{callee}}}}

% Listing style
\usepackage{listings}
\lstset{
  basicstyle=\ttfamily\footnotesize,
  breaklines=true,
  frame=single,
  numbers=left,
  numberstyle=\tiny,
  captionpos=b,
}
 for arXiv

\ifanonymous
  \newcommand{\toolname}{{\sc BoundaryFlow}\xspace}
  \newcommand{\anonrepo}{\url{https://anonymous.4open.science/r/XXXX}}
  \newcommand{\authorblock}{%
    \author{Anonymous Submission \\ Paper ID: TBD}%
  }
\else
  \newcommand{\toolname}{{\sc kcov-dataflow}\xspace}
  \newcommand{\anonrepo}{\url{https://github.com/yskzalloc/kcov-dataflow}}
  \newcommand{\authorblock}{%
    \author{Yunseong Kim \\
      Ericsson Software Technology \\
      \texttt{yunseong.kim@est.tech}}%
    % \author{Anonymous Submission}%
  }
\fi

% Technical shorthands
\newcommand{\sancov}{\texttt{SanitizerCoverage}\xspace}
\newcommand{\tracepc}{\texttt{trace-pc}\xspace}
\newcommand{\traceargs}{\texttt{trace-args}\xspace}
\newcommand{\traceret}{\texttt{trace-ret}\xspace}
\newcommand{\kcov}{\texttt{KCOV}\xspace}
\newcommand{\kasan}{\texttt{KASAN}\xspace}
\newcommand{\syzkaller}{\texttt{syzkaller}\xspace}
\newcommand{\drgn}{\texttt{drgn}\xspace}
\newcommand{\virtme}{\texttt{virtme-ng}\xspace}

% Data-flow tuple notation
\newcommand{\dftuple}[4]{\ensuremath{\langle #1, #2, #3, #4 \rangle}}
\newcommand{\callerid}{\ensuremath{\mathit{PC}_{\mathrm{caller}}}}
\newcommand{\calleeid}{\ensuremath{\mathit{PC}_{\mathrm{callee}}}}

% Listing style
\usepackage{listings}
\lstset{
  basicstyle=\ttfamily\footnotesize,
  breaklines=true,
  frame=single,
  numbers=left,
  numberstyle=\tiny,
  captionpos=b,
}
 for NDSS,
%                       set \anonymousfalse before % shared/macros.tex — Identity-masked macros for both arXiv and NDSS versions
% Toggle anonymization: set \anonymoustrue before % shared/macros.tex — Identity-masked macros for both arXiv and NDSS versions
% Toggle anonymization: set \anonymoustrue before \input{macros} for NDSS,
%                       set \anonymousfalse before \input{macros} for arXiv

\ifanonymous
  \newcommand{\toolname}{{\sc BoundaryFlow}\xspace}
  \newcommand{\anonrepo}{\url{https://anonymous.4open.science/r/XXXX}}
  \newcommand{\authorblock}{%
    \author{Anonymous Submission \\ Paper ID: TBD}%
  }
\else
  \newcommand{\toolname}{{\sc kcov-dataflow}\xspace}
  \newcommand{\anonrepo}{\url{https://github.com/yskzalloc/kcov-dataflow}}
  \newcommand{\authorblock}{%
    \author{Yunseong Kim \\
      Ericsson Software Technology \\
      \texttt{yunseong.kim@est.tech}}%
    % \author{Anonymous Submission}%
  }
\fi

% Technical shorthands
\newcommand{\sancov}{\texttt{SanitizerCoverage}\xspace}
\newcommand{\tracepc}{\texttt{trace-pc}\xspace}
\newcommand{\traceargs}{\texttt{trace-args}\xspace}
\newcommand{\traceret}{\texttt{trace-ret}\xspace}
\newcommand{\kcov}{\texttt{KCOV}\xspace}
\newcommand{\kasan}{\texttt{KASAN}\xspace}
\newcommand{\syzkaller}{\texttt{syzkaller}\xspace}
\newcommand{\drgn}{\texttt{drgn}\xspace}
\newcommand{\virtme}{\texttt{virtme-ng}\xspace}

% Data-flow tuple notation
\newcommand{\dftuple}[4]{\ensuremath{\langle #1, #2, #3, #4 \rangle}}
\newcommand{\callerid}{\ensuremath{\mathit{PC}_{\mathrm{caller}}}}
\newcommand{\calleeid}{\ensuremath{\mathit{PC}_{\mathrm{callee}}}}

% Listing style
\usepackage{listings}
\lstset{
  basicstyle=\ttfamily\footnotesize,
  breaklines=true,
  frame=single,
  numbers=left,
  numberstyle=\tiny,
  captionpos=b,
}
 for NDSS,
%                       set \anonymousfalse before % shared/macros.tex — Identity-masked macros for both arXiv and NDSS versions
% Toggle anonymization: set \anonymoustrue before \input{macros} for NDSS,
%                       set \anonymousfalse before \input{macros} for arXiv

\ifanonymous
  \newcommand{\toolname}{{\sc BoundaryFlow}\xspace}
  \newcommand{\anonrepo}{\url{https://anonymous.4open.science/r/XXXX}}
  \newcommand{\authorblock}{%
    \author{Anonymous Submission \\ Paper ID: TBD}%
  }
\else
  \newcommand{\toolname}{{\sc kcov-dataflow}\xspace}
  \newcommand{\anonrepo}{\url{https://github.com/yskzalloc/kcov-dataflow}}
  \newcommand{\authorblock}{%
    \author{Yunseong Kim \\
      Ericsson Software Technology \\
      \texttt{yunseong.kim@est.tech}}%
    % \author{Anonymous Submission}%
  }
\fi

% Technical shorthands
\newcommand{\sancov}{\texttt{SanitizerCoverage}\xspace}
\newcommand{\tracepc}{\texttt{trace-pc}\xspace}
\newcommand{\traceargs}{\texttt{trace-args}\xspace}
\newcommand{\traceret}{\texttt{trace-ret}\xspace}
\newcommand{\kcov}{\texttt{KCOV}\xspace}
\newcommand{\kasan}{\texttt{KASAN}\xspace}
\newcommand{\syzkaller}{\texttt{syzkaller}\xspace}
\newcommand{\drgn}{\texttt{drgn}\xspace}
\newcommand{\virtme}{\texttt{virtme-ng}\xspace}

% Data-flow tuple notation
\newcommand{\dftuple}[4]{\ensuremath{\langle #1, #2, #3, #4 \rangle}}
\newcommand{\callerid}{\ensuremath{\mathit{PC}_{\mathrm{caller}}}}
\newcommand{\calleeid}{\ensuremath{\mathit{PC}_{\mathrm{callee}}}}

% Listing style
\usepackage{listings}
\lstset{
  basicstyle=\ttfamily\footnotesize,
  breaklines=true,
  frame=single,
  numbers=left,
  numberstyle=\tiny,
  captionpos=b,
}
 for arXiv

\ifanonymous
  \newcommand{\toolname}{{\sc BoundaryFlow}\xspace}
  \newcommand{\anonrepo}{\url{https://anonymous.4open.science/r/XXXX}}
  \newcommand{\authorblock}{%
    \author{Anonymous Submission \\ Paper ID: TBD}%
  }
\else
  \newcommand{\toolname}{{\sc kcov-dataflow}\xspace}
  \newcommand{\anonrepo}{\url{https://github.com/yskzalloc/kcov-dataflow}}
  \newcommand{\authorblock}{%
    \author{Yunseong Kim \\
      Ericsson Software Technology \\
      \texttt{yunseong.kim@est.tech}}%
    % \author{Anonymous Submission}%
  }
\fi

% Technical shorthands
\newcommand{\sancov}{\texttt{SanitizerCoverage}\xspace}
\newcommand{\tracepc}{\texttt{trace-pc}\xspace}
\newcommand{\traceargs}{\texttt{trace-args}\xspace}
\newcommand{\traceret}{\texttt{trace-ret}\xspace}
\newcommand{\kcov}{\texttt{KCOV}\xspace}
\newcommand{\kasan}{\texttt{KASAN}\xspace}
\newcommand{\syzkaller}{\texttt{syzkaller}\xspace}
\newcommand{\drgn}{\texttt{drgn}\xspace}
\newcommand{\virtme}{\texttt{virtme-ng}\xspace}

% Data-flow tuple notation
\newcommand{\dftuple}[4]{\ensuremath{\langle #1, #2, #3, #4 \rangle}}
\newcommand{\callerid}{\ensuremath{\mathit{PC}_{\mathrm{caller}}}}
\newcommand{\calleeid}{\ensuremath{\mathit{PC}_{\mathrm{callee}}}}

% Listing style
\usepackage{listings}
\lstset{
  basicstyle=\ttfamily\footnotesize,
  breaklines=true,
  frame=single,
  numbers=left,
  numberstyle=\tiny,
  captionpos=b,
}
 for arXiv

\ifanonymous
  \newcommand{\toolname}{{\sc BoundaryFlow}\xspace}
  \newcommand{\anonrepo}{\url{https://anonymous.4open.science/r/XXXX}}
  \newcommand{\authorblock}{%
    \author{Anonymous Submission \\ Paper ID: TBD}%
  }
\else
  \newcommand{\toolname}{{\sc kcov-dataflow}\xspace}
  \newcommand{\anonrepo}{\url{https://github.com/yskzalloc/kcov-dataflow}}
  \newcommand{\authorblock}{%
    \author{Yunseong Kim \\
      Ericsson Software Technology \\
      \texttt{yunseong.kim@est.tech}}%
    % \author{Anonymous Submission}%
  }
\fi

% Technical shorthands
\newcommand{\sancov}{\texttt{SanitizerCoverage}\xspace}
\newcommand{\tracepc}{\texttt{trace-pc}\xspace}
\newcommand{\traceargs}{\texttt{trace-args}\xspace}
\newcommand{\traceret}{\texttt{trace-ret}\xspace}
\newcommand{\kcov}{\texttt{KCOV}\xspace}
\newcommand{\kasan}{\texttt{KASAN}\xspace}
\newcommand{\syzkaller}{\texttt{syzkaller}\xspace}
\newcommand{\drgn}{\texttt{drgn}\xspace}
\newcommand{\virtme}{\texttt{virtme-ng}\xspace}

% Data-flow tuple notation
\newcommand{\dftuple}[4]{\ensuremath{\langle #1, #2, #3, #4 \rangle}}
\newcommand{\callerid}{\ensuremath{\mathit{PC}_{\mathrm{caller}}}}
\newcommand{\calleeid}{\ensuremath{\mathit{PC}_{\mathrm{callee}}}}

% Listing style
\usepackage{listings}
\lstset{
  basicstyle=\ttfamily\footnotesize,
  breaklines=true,
  frame=single,
  numbers=left,
  numberstyle=\tiny,
  captionpos=b,
}


% Header
\pagestyle{fancy}
\thispagestyle{empty}
\renewcommand{\shorttitle}{Per-Task Data-Flow Extraction at Kernel Function Boundaries}
\renewcommand{\headeright}{A Preprint}

%% Title
\title{Beyond Edge Coverage: Per-Task Data-Flow Extraction at Kernel Function Boundaries via LLVM
}

\authorblock
\date{}

\begin{document}
\maketitle

\begin{abstract}
Coverage-guided kernel fuzzers such as syzkaller rely on edge coverage
(\texttt{trace-pc}) as their sole feedback signal. This context-blind approach
cannot distinguish execution paths that differ only in argument values---for
example, two invocations of \texttt{copy\_from\_user()} with different size
parameters hit identical basic blocks yet have vastly different security
implications. I present \toolname{}, an LLVM-based instrumentation framework
that extends Linux KCOV with data-flow extraction of function arguments and
return values. A compiler pass emits lightweight callbacks capturing structured
tuples of program counter, argument metadata, and field values at function
entry and return. Composite types are automatically decomposed via DWARF
\texttt{DICompositeType} metadata with zero source annotation. A lock-free
per-task ring buffer delivers records to user space with no interference to
existing KCOV or syzkaller infrastructure. I demonstrate dual utility:
(1)~fuzzers gain state-aware feedback for mutation guidance into
value-dependent state transitions, and (2)~security analysts obtain
deterministic argument records for root-cause analysis without
\texttt{printk} or kprobe overhead. Two Rust instrumentation paths are
provided: a post-compilation pipeline requiring no rustc modification, and
native instrumentation via rustc built against the custom LLVM---both the
only runtime methods for capturing Rust function arguments given that
\texttt{drgn}/vmcore fails under \texttt{-O2} DWARF elision.
\end{abstract}

\keywords{Kernel Fuzzing \and Data-Flow Extraction \and LLVM Instrumentation \and KCOV \and Rust for Linux}


\section{Introduction}\label{sec:intro}

The Linux kernel exceeds 40 million lines of code, integrates Rust(~397k Lines) as a second
systems language, and implements complex locking hierarchies
(\texttt{rwsem}, \texttt{mutex}, RCU) that create deep state spaces unreachable
by coverage-guided fuzzing alone. syzkaller~\cite{vyukov2015syzkaller} has
discovered over 5{,}000 kernel bugs in the linux kernel, yet plateaus on stateful
subsystems such as \texttt{ksmbd}, \texttt{io\_uring}, and \texttt{binder}
where security-critical behavior depends on runtime argument values rather than
control-flow topology.

The fundamental limitation is a \emph{semantic gap}: KCOV's \tracepc reports
\emph{which} basic blocks executed but not \emph{with what values}. Two
executions of \texttt{vfs\_open()} with different access-mode flags traverse
identical edges yet have entirely different security implications. Fuzzers
cannot distinguish these without data-flow context, leading to coverage
saturation on value-dependent state transitions.

Prior approaches each fall short for kernel-scale deployment. Dynamic tracing
via \texttt{ftrace}/\texttt{kprobes} imposes per-function overhead unsuitable
for continuous fuzzing. IJON~\cite{aschermann2020ijon} requires manual
annotation---infeasible for 30M~LOC. Context-sensitive coverage
(Angora~\cite{chen2018angora}, CollAFL~\cite{gan2018collafl}) clones call sites
but explodes on polymorphic dispatch (\texttt{file\_operations} alone has 20+
function pointers). DataFlowSanitizer provides byte-level taint tracking but
imposes 10--100$\times$ overhead, rendering it unusable for kernel fuzzing.

I observe that \emph{function boundaries} are natural chokepoints where
data-flow context is both cheap to capture and highly informative. Arguments
encode the \emph{intent} of a call; return values encode the \emph{outcome}.
Capturing these at compile time with a lock-free kernel transport gives both
fuzzers and analysts a new dimension of observability without source annotation.

\toolname{} is an automated LLVM IR instrumentation framework that captures the
data-flow tuple $\dftuple{\callerid}{\mathit{arg\_idx} \times
\mathit{arg\_size}}{\mathit{ptr}}{\mathit{offsets}[]}$ at every function
boundary. It operates through three layers:
\begin{enumerate}[nosep]
  \item A compiler pass that emits callbacks with struct field expansion via
        DWARF metadata.
  \item A kernel-resident lock-free per-task ring buffer delivered through a
        separate \texttt{/sys/kernel/debug/kcov\_dataflow} dedicated debugfs interface.
  \item A user-space consumer interface for both fuzzer executors and human
        analysts.
\end{enumerate}

My contributions are:
\begin{itemize}[nosep]
  \item An LLVM \sancov extension\footnote{\url{https://github.com/llvm/llvm-project/pull/201410}} that captures function arguments and return
        values with automatic struct field expansion via \texttt{DICompositeType}
        metadata, requiring zero source annotation.
  \item A lock-free, per-task kernel transport mechanism that operates safely in
        process context with explicit interrupt rejection, without interfering
        with existing KCOV or syzkaller infrastructure.
  \item A post-compilation pipeline enabling Rust kernel module instrumentation
        without modifying the Rust compiler toolchain, and a native path via
        rustc built against the custom LLVM.
  \item Empirical demonstration on five vulnerability classes showing that
        boundary data-flow context provides information unavailable through any
        combination of existing tools (coverage + sanitizers + post-mortem
        debuggers).
\end{itemize}


\section{Background \& Motivation}\label{sec:background}

\subsection{KCOV and Coverage-Guided Kernel Fuzzing}

KCOV~\cite{vyukov2016kcov} instruments the kernel at compile time via
\sancov's \tracepc callback. Each task owns a kernel-allocated buffer,
memory-mapped into user space, into which the kernel writes the program
counter of every executed edge. User-space fuzzers (primarily syzkaller)
hash these PCs into a coverage signal, retaining inputs that discover
novel edges in the corpus.

This architecture has proven remarkably effective: syzkaller continuously
discovers 100+ bugs per month across upstream kernels. However, coverage
\emph{saturates} on stateful subsystems. After initial exploration, new edges
become rare, and the fuzzer degenerates into random mutation without meaningful
guidance.

\subsection{The Data-Flow Blind Spot}

\begin{figure}[t]
\centering
\begin{tikzpicture}[
  node distance=0.3cm,
  box/.style={rectangle, draw, thick, minimum width=2.2cm, minimum height=0.6cm, font=\scriptsize\bfseries},
  msg/.style={-{Stealth[length=2mm]}, thick},
  lbl/.style={font=\scriptsize, midway, above},
  note/.style={font=\scriptsize\itshape, text=gray}
]
\node[box] (client) {Client (Fuzzer)};
\node[box, right=5cm of client] (server) {ksmbd (Kernel)};

% Vertical lifelines
\draw[dashed, gray] (client.south) -- ++(0,-5.5);
\draw[dashed, gray] (server.south) -- ++(0,-5.5);

% Messages
\draw[msg] ([yshift=-0.8cm]client.south) -- node[lbl] {SMB2\_NEGOTIATE} ([yshift=-0.8cm]server.south);
\draw[msg] ([yshift=-1.4cm]server.south) -- node[lbl] {dialect OK} ([yshift=-1.4cm]client.south);

\draw[msg] ([yshift=-2.0cm]client.south) -- node[lbl] {SMB2\_SESSION\_SETUP} ([yshift=-2.0cm]server.south);
\draw[msg] ([yshift=-2.6cm]server.south) -- node[lbl] {authenticated} ([yshift=-2.6cm]client.south);

\draw[msg] ([yshift=-3.2cm]client.south) -- node[lbl] {SMB2\_TREE\_CONNECT} ([yshift=-3.2cm]server.south);
\draw[msg] ([yshift=-3.8cm]server.south) -- node[lbl] {share bound} ([yshift=-3.8cm]client.south);

\draw[msg, thick, red!70!black] ([yshift=-4.6cm]client.south) -- node[lbl, text=red!70!black] {SMB2\_CREATE(\colorbox{yellow!30}{DesiredAccess})} ([yshift=-4.6cm]server.south);
\draw[msg] ([yshift=-5.2cm]server.south) -- node[lbl] {file handle} ([yshift=-5.2cm]client.south);

% Annotations
\node[note, anchor=west] at ([xshift=0.5cm, yshift=-4.6cm]server.south) {$\to$ vfs\_open(flags)};
\node[note, anchor=east] at ([xshift=-0.3cm, yshift=-1.1cm]client.south) {coverage saturates here};
\draw[decorate, decoration={brace, amplitude=3pt, mirror, raise=2pt}, gray]
  ([xshift=0.2cm, yshift=-0.6cm]client.south) -- ([xshift=0.2cm, yshift=-3.9cm]client.south)
  node[midway, left=5pt, note] {same edges};
\end{tikzpicture}
\caption{SMB2 protocol stages required before reaching \texttt{vfs\_open()}.
Edge coverage saturates after discovering the handshake sequence; the
security-relevant \texttt{DesiredAccess} field inside SMB2\_CREATE is
invisible to control-flow feedback.}\label{fig:ksmbd}
\end{figure}

Consider the \texttt{ksmbd} SMB2\_CREATE handler. Before a file operation
reaches the VFS layer, ksmbd enforces a multi-stage protocol handshake
between the client (fuzzer) and the in-kernel SMB server
(Figure~\ref{fig:ksmbd}). A fuzzer must produce valid negotiate, authenticate,
and tree-connect messages before it can exercise the file-create path. Once
coverage of this multi-stage sequence saturates, all subsequent CREATE requests
traverse identical edges regardless of \texttt{DesiredAccess}---the fuzzer
cannot distinguish \texttt{FILE\_READ\_DATA} from
\texttt{FILE\_WRITE\_DATA|FILE\_DELETE\_ON\_CLOSE}. With \toolname{}, the
\texttt{DesiredAccess} field is captured at the \texttt{smb2\_open()} boundary,
enabling value-aware mutation guidance at deep protocol stages.

A second example arises at Rust-to-C FFI boundaries. A Rust wrapper calls a C
function with a struct pointer whose fields determine kernel behavior. Edge
coverage sees one call site; the struct contents---which control whether a
privilege check is bypassed---are invisible.

\begin{figure}[t]
\centering
\begin{tikzpicture}[
  node distance=0.4cm and 1.2cm,
  block/.style={rectangle, draw, rounded corners, minimum width=2.2cm, minimum height=0.7cm, font=\footnotesize, align=center},
  arrow/.style={-{Stealth[length=2mm]}, thick},
  label/.style={font=\scriptsize\itshape, text=gray}
]
% Left side: Edge Coverage
\node[font=\footnotesize\bfseries] (t1) {Edge Coverage (KCOV)};
\node[block, below=0.3cm of t1] (a1) {syscall()};
\node[block, below=of a1] (b1) {vfs\_open()};
\node[block, below=of b1] (c1) {do\_open()};
\draw[arrow] (a1) -- (b1);
\draw[arrow] (b1) -- (c1);
\node[label, right=0.1cm of b1] {flags=?};
\node[label, below=0.1cm of c1] {Sees: A$\to$B$\to$C};

% Right side: Data-Flow Extraction
\node[font=\footnotesize\bfseries, right=3.5cm of t1] (t2) {+ Data-Flow Extraction};
\node[block, below=0.3cm of t2] (a2) {syscall()};
\node[block, below=of a2, fill=green!10] (b2) {vfs\_open()};
\node[block, below=of b2] (c2) {do\_open()};
\draw[arrow] (a2) -- (b2);
\draw[arrow] (b2) -- (c2);
\node[label, right=0.1cm of b2, text=black] {flags=O\_WRONLY|O\_TRUNC};
\node[label, below=0.1cm of c2] {Sees: A$\to$B(flags=0x241)$\to$C};
\end{tikzpicture}
\caption{Edge coverage reports identical paths for different argument values.
Data-flow extraction reveals the security-relevant difference.}\label{fig:blind-spot}
\end{figure}

\begin{figure}[t]
\centering
\begin{tikzpicture}[
  node distance=0.15cm,
  rec/.style={font=\ttfamily\scriptsize, anchor=west},
  entry/.style={rec, text=black!80},
  ret/.style={rec, text=black!60},
  hl/.style={fill=yellow!30, rounded corners=1pt, inner sep=1pt},
]
\node[entry] (l1) at (0,0) {do\_sys\_openat2(0xffffff9c, \colorbox{yellow!30}{"/etc/shadow"}, 0x241)};
\node[entry] (l2) at (0.4,-0.4) {build\_open\_flags(\colorbox{yellow!30}{O\_WRONLY|O\_TRUNC})};
\node[ret]   (l3) at (0.4,-0.8) {0x0 = build\_open\_flags()};
\node[entry] (l4) at (0.4,-1.2) {getname\_flags(\colorbox{yellow!30}{"/etc/shadow"}, 0x0)};
\node[ret]   (l5) at (0.4,-1.6) {0xffff8880... = getname\_flags()};
\node[entry] (l6) at (0.4,-2.0) {vfs\_open(path, \colorbox{yellow!30}{\{.flags=0x241, .mode=0x1b6\}})};
\node[entry] (l7) at (0.8,-2.4) {do\_dentry\_open(file, inode, 0x0)};
\node[entry] (l8) at (1.2,-2.8) {security\_file\_open(\colorbox{yellow!30}{file, 0x241})};
\node[ret]   (l9) at (1.2,-3.2) {0x0 = security\_file\_open()};
\node[ret]   (la) at (0.8,-3.6) {0x0 = do\_dentry\_open()};
\node[ret]   (lb) at (0.4,-4.0) {\colorbox{green!20}{0x4} = vfs\_open()};
\node[ret]   (lc) at (0,-4.4) {\colorbox{green!20}{0x4} = do\_sys\_openat2()};

% Annotation
\draw[decorate, decoration={brace, amplitude=4pt, mirror, raise=3pt}]
  ([xshift=7.5cm]l1.north east) -- ([xshift=7.5cm]lc.south east)
  node[midway, right=7pt, font=\scriptsize\itshape, align=left] {Captured\\per-task\\at runtime};
\end{tikzpicture}
\caption{Actual \toolname{} output for \texttt{open("/etc/shadow", O\_WRONLY|O\_TRUNC)}.
Each line is a captured record showing function arguments (ENTRY) and return values (RET)
with automatic struct field expansion. Indentation reflects call depth.}\label{fig:calltree}
\end{figure}

\subsection{Threat Model and Assumptions}

I assume a standard kernel fuzzing threat model: an unprivileged user-space
process with access to the syscall interface attempts to trigger memory
corruption in kernel subsystems. KASAN/KMSAN detect the corruption once
triggered; my tool provides the \emph{path to triggering} it and the
\emph{data for triaging} it. I assume the kernel is compiled with debug
information (\texttt{-g}) for full struct field expansion, though the system
degrades gracefully to scalar-only capture without it.


\section{System Design \& Architecture}\label{sec:design}

\subsection{Overview}

\toolname{} operates as a three-layer pipeline: (1)~compile-time
instrumentation inserts data-flow callbacks into kernel object files,
(2)~a runtime transport layer delivers records to user-space via a lock-free
per-task ring buffer, and (3)~user-space consumers (fuzzer executors or analyst
tools) read the buffer via \texttt{mmap()}. The system is activated by the
compiler flags \texttt{-fsanitize-coverage=trace-args,trace-ret} and
requires no source-level annotation.

\begin{figure}[t]
\centering
\begin{tikzpicture}[
  node distance=0.5cm and 0.8cm,
  layer/.style={rectangle, draw, thick, rounded corners, minimum width=9cm, minimum height=1.6cm, font=\footnotesize, align=center},
  comp/.style={rectangle, draw, fill=blue!8, minimum width=2.5cm, minimum height=0.6cm, font=\scriptsize, align=center},
  arrow/.style={-{Stealth[length=2mm]}, thick},
  label/.style={font=\scriptsize\itshape}
]
% Layer 1: Compile time
\node[layer, fill=yellow!10] (l1) {};
\node[font=\scriptsize\bfseries, anchor=north west] at ([xshift=3pt,yshift=-3pt]l1.north west) {Compile Time};
\node[comp] at ([yshift=-0.2cm]l1.center) (pass) {LLVM SanCov Pass\\(trace-args/ret)};

% Layer 2: Kernel runtime
\node[layer, fill=orange!10, below=0.5cm of l1] (l2) {};
\node[font=\scriptsize\bfseries, anchor=north west] at ([xshift=3pt,yshift=-3pt]l2.north west) {Kernel Runtime};
\node[comp, fill=orange!15] at ([xshift=-2cm,yshift=-0.2cm]l2.center) (cb) {\_\_sanitizer\_cov\\\_trace\_args/ret};
\node[comp, fill=orange!15] at ([xshift=2cm,yshift=-0.2cm]l2.center) (buf) {Per-Task\\Ring Buffer};

% Layer 3: User space
\node[layer, fill=green!10, below=0.5cm of l2] (l3) {};
\node[font=\scriptsize\bfseries, anchor=north west] at ([xshift=3pt,yshift=-3pt]l3.north west) {User Space};
\node[comp, fill=green!15] at ([yshift=-0.2cm]l3.center) (user) {mmap() consumer\\(fuzzer / analyst)};

% Arrows
\draw[arrow] (pass) -- node[left, label, pos=0.35] {inserts callbacks} (cb);
\draw[arrow] (cb) -- node[above, label] {atomic write} (buf);
\draw[arrow] (buf) -- node[left, label, pos=0.35] {shared memory} (user);
\end{tikzpicture}
\caption{Three-layer architecture: compile-time instrumentation, kernel-resident
lock-free buffer, and user-space consumption via \texttt{mmap()}.}\label{fig:arch}
\end{figure}

\subsection{LLVM IR Boundary Instrumentation}\label{sec:llvm-pass}

The instrumentation pass resides in LLVM's existing \sancov framework
(\texttt{SanitizerCoverage.cpp}). For each function~$F$ with a
\texttt{DISubprogram} debug info attachment:

\paragraph{Argument Capture.} At function entry, for each parameter~$i$:
\begin{enumerate}[nosep]
  \item Resolve the parameter's \texttt{DILocalVariable} $\to$ \texttt{DIType}.
  \item If the type is a \texttt{DICompositeType} (struct/union): walk its
        \texttt{DIDerivedType} members to extract field byte-offsets and sizes.
        Compute an FNV-1a hash of the type name for deduplication.
  \item Emit a compile-time constant global array:
        \texttt{[hash, off$_0$, sz$_0$, off$_1$, sz$_1$, \ldots, off$_N$, sz$_N$]}.
  \item Insert a call to \texttt{\_\_sanitizer\_cov\_trace\_args(pc, arg\_idx,
        arg\_size, ptr, offsets\_ptr, num\_fields)} where \texttt{offsets\_ptr}
        points past the hash element (GEP index~1).
\end{enumerate}
For scalar arguments, \texttt{offsets\_ptr} is \texttt{null} and
\texttt{num\_fields} is~0. Pointer arguments are passed directly; scalars are
spilled to an \texttt{alloca} so the kernel callback receives a uniform pointer
interface.

\paragraph{Return Value Capture.} Using LLVM's \texttt{EscapeEnumerator}, the
pass locates every \texttt{ReturnInst} in~$F$ and inserts a call to
\texttt{\_\_sanitizer\_cov\_trace\_ret(pc, ret\_size, ptr, offsets\_ptr,
num\_fields)} before each return.

\paragraph{Edge Cases.} Functions without debug info, variadic functions, and
\texttt{naked} functions are skipped gracefully. The pass is controlled by two
\texttt{cl::opt} flags: \texttt{-sanitizer-coverage-trace-args} and
\texttt{-sanitizer-coverage-trace-ret}.

\subsection{Lock-Free Per-Task Ring Buffer}\label{sec:ringbuf}

The kernel backend adds three fields to \texttt{task\_struct}:
\texttt{kcov\_df\_area} (pointer to \texttt{u64} buffer),
\texttt{kcov\_df\_size} (buffer capacity), and \texttt{kcov\_df\_seq}
(monotonic sequence counter).

\paragraph{Write Path.} The callback implementation uses
\texttt{READ\_ONCE}/\texttt{WRITE\_ONCE} to reserve a contiguous slot in
the buffer (portable to ARM64 without architecture-specific atomics).
If the reserved offset exceeds
\texttt{kcov\_df\_size}, the event is silently dropped---no blocking, no
allocation, no spinlock. All pointer dereferences use
\texttt{copy\_from\_kernel\_nofault()} to safely handle unmapped or poisoned
memory (e.g., KASAN red zones, freed slabs). Values in the
\texttt{IS\_ERR\_VALUE} range ($[-4095, -1]$) are detected and skipped.

\paragraph{Record Format.} Each record is a variable-length sequence of
\texttt{u64} words starting with a 3-word header: \texttt{[type\_and\_seq | pc | meta | field\_val$_0$ | \ldots | field\_val$_N$]} (or a single scalar value if \texttt{num\_fields} is 0). Here, \texttt{type\_and\_seq} packs the type marker (\texttt{0xE0000000} for entry, \texttt{0xF0000000} for return) and the 24-bit sequence number. The \texttt{meta} word packs the 8-bit argument index (for entry, 0 for return), 8-bit argument or return size in bytes, and the lower 48 bits of the raw pointer value.

\begin{figure}[t]
\centering
\begin{tikzpicture}[
  cell/.style={rectangle, draw, minimum width=1.4cm, minimum height=0.5cm, font=\tiny, align=center},
  hdr/.style={cell, fill=blue!12},
  dat/.style={cell, fill=green!12},
  cnt/.style={cell, fill=gray!15},
]
\node[cnt] (c0) {area[0]\\count};
\node[hdr, right=0pt of c0] (h1) {type|seq};
\node[hdr, right=0pt of h1] (h2) {PC};
\node[hdr, right=0pt of h2] (h3) {meta};
\node[dat, right=0pt of h3] (d1) {field$_0$};
\node[dat, right=0pt of d1] (d2) {field$_1$};
\node[cell, right=0pt of d2, draw=none] (dots) {\ldots};

\node[font=\tiny, below=0.2cm of c0] {u64};
\node[font=\tiny\itshape, above=0.15cm of h1] {3-word header};
\node[font=\tiny\itshape, above=0.15cm of d1] {N fields};

\draw[decorate, decoration={brace, amplitude=3pt, raise=2pt}]
  (h1.north west) -- (h3.north east);
\draw[decorate, decoration={brace, amplitude=3pt, raise=2pt}]
  (d1.north west) -- (d2.north east);
\end{tikzpicture}
\caption{Ring buffer layout: area[0] holds the atomic write counter;
each record has a 3-word header followed by field values.}\label{fig:record}
\end{figure}

\paragraph{Safety.} Callbacks are marked \texttt{notrace},
\texttt{\_\_no\_sanitize\_coverage}, and \texttt{noinline} to prevent recursion.
No \texttt{printk} or allocation occurs in the data path. An
\texttt{in\_task()} guard at the top of the write path rejects all calls from
non-process contexts (softirq, hardirq, and NMI), eliminating the risk of
reentrant corruption when a hardware or software interrupt preempts a task
that is already inside the callback.

\subsection{Separate Device Interface}\label{sec:device}

\toolname{} exposes \texttt{/sys/kernel/debug/kcov\_dataflow} as a device
\emph{completely independent} from the exsiting \texttt{/sys/kernel/debug/kcov}.
This ensures zero interference with syzkaller's existing KCOV usage.

The ioctl namespace uses magic~\texttt{'d'}:
\begin{itemize}[nosep]
  \item \texttt{KCOV\_DF\_INIT\_TRACK} (\texttt{0x80086401}): allocate buffer
  \item \texttt{KCOV\_DF\_ENABLE} (\texttt{0x6464}): start recording
  \item \texttt{KCOV\_DF\_DISABLE} (\texttt{0x6465}): stop recording
\end{itemize}
User-space maps the buffer via \texttt{mmap()} for zero-copy reads.

\subsection{Rust Module Support}\label{sec:rust}

The Rust compiler (\texttt{rustc}) bundles its own LLVM (v22 for rustc~1.95)
which lacks my custom flags. Two instrumentation paths exist:

\paragraph{Path A: Post-compilation pipeline (no rustc modification).}
For environments using an unmodified rustc:
\begin{enumerate}[nosep]
  \item \texttt{rustc --emit=llvm-ir}: generate LLVM IR text with debug info
  \item \texttt{opt -passes=sancov-module -sanitizer-coverage-trace-args
        -sanitizer-coverage-trace-ret}: instrument with my pass
  \item \texttt{llc -filetype=obj}: compile to object code
  \item Standard module linking produces the instrumented \texttt{.ko}
\end{enumerate}
This works because LLVM IR is language-agnostic: my pass operates identically
on Rust-generated IR. The text format (\texttt{.ll}) avoids bitcode version
incompatibilities between LLVM~22 and~23. Struct field expansion works on Rust
\texttt{\#[repr(C)]} types whose \texttt{DICompositeType} metadata is preserved
by \texttt{rustc -g}.

\begin{figure}[t]
\centering
\begin{tikzpicture}[
  node distance=0.3cm and 0.4cm,
  pstep/.style={rectangle, draw, fill=blue!8, minimum width=1.8cm, minimum height=0.6cm, font=\scriptsize, align=center},
  arrow/.style={-{Stealth[length=2mm]}, thick},
  label/.style={font=\tiny\itshape, text=gray}
]
\node[pstep] (rustc) {rustc\\--emit=llvm-ir};
\node[pstep, right=0.6cm of rustc] (opt) {opt\\(my pass)};
\node[pstep, right=0.6cm of opt] (llc) {llc\\-filetype=obj};
\node[pstep, right=0.6cm of llc] (ld) {ld\\module.ko};

\draw[arrow] (rustc) -- node[above, label] {.ll} (opt);
\draw[arrow] (opt) -- node[above, label] {.ll} (llc);
\draw[arrow] (llc) -- node[above, label] {.o} (ld);

\node[label, below=0.15cm of rustc] {Rust LLVM 22};
\node[label, below=0.15cm of opt] {Custom LLVM 23};
\node[label, below=0.15cm of llc] {Custom LLVM 23};
\end{tikzpicture}
\caption{Rust post-compilation pipeline: text-format IR avoids bitcode version
incompatibility between rustc's bundled LLVM and my custom pass.}\label{fig:rust-pipeline}
\end{figure}

\paragraph{Path B: Native instrumentation (custom rustc).}
Building \texttt{rustc} against the custom LLVM~23 (via
\texttt{llvm-config} in \texttt{config.toml}) gives rustc native support for
\texttt{-Cllvm-args=-sanitizer-coverage-trace-args/ret}. The kernel build
system then passes these flags transparently via
\texttt{RUSTFLAGS\_KCOV\_DATAFLOW}---no post-compilation step required.
\textbf{This is the required path for tracing Rust kernel modules}, as the
kernel's Kbuild compiles \texttt{core}, \texttt{alloc}, and \texttt{kernel}
crates internally and does not expose intermediate IR for external
instrumentation.


\subsection{Pointer Dereference and Dynamic Struct Capture}\label{sec:deref}

A key design question is whether \toolname{} can capture struct fields from
dynamically allocated (heap) memory, not just stack-local values. The answer is
yes: the kernel callback performs \emph{one level of pointer dereference} at
runtime.

\paragraph{Mechanism.} When a function parameter is a pointer to a struct
(e.g., \texttt{struct session\_data *sd} where \texttt{sd} was
\texttt{kmalloc}'d), the LLVM pass emits the pointer value and the compile-time
offset array. At runtime, the callback computes each field address as
$\mathit{field\_addr} = \mathit{ptr} + \mathit{offset}[i]$ and reads via
\texttt{copy\_from\_kernel\_nofault()}. This works regardless of whether the
struct resides on the heap, stack, or in static memory.

\paragraph{Safety for edge cases:}
\begin{itemize}[nosep]
  \item \textbf{Freed slab (KASAN poison):} Safely reads poison bytes
        (\texttt{0xfd}/\texttt{0xff})---revealing the object was freed.
  \item \textbf{Unmapped page:} \texttt{copy\_from\_kernel\_nofault()} returns
        error; field stored as~0.
  \item \textbf{ERR\_PTR:} Detected via \texttt{IS\_ERR\_VALUE} range check;
        pointer not dereferenced.
  \item \textbf{NULL:} Detected and skipped (no crash).
  \item \textbf{Nested structs:} Flattened---\texttt{DICompositeType} includes
        nested fields with correct absolute offsets.
\end{itemize}

\paragraph{Design choice: shallow capture.} I perform exactly one level of
dereference (pointer $\to$ struct fields). I do \emph{not} chase pointers
inside the struct (e.g., a \texttt{char *buf} field stores the pointer value,
not the string contents). This bounds the cost to $O(n)$ where $n$ is the
number of struct fields, and avoids recursive faults on invalid nested pointers.
The pointer value itself is informative: NULL indicates freed/uninitialized,
\texttt{0xffff888...} indicates valid kernel heap, and KASAN poison patterns
indicate use-after-free.

\section{Implementation}\label{sec:impl}

\subsection{LLVM Pass (270~LOC)}

The pass extends \texttt{SanitizerCoverage.cpp} with two new methods:

\texttt{InjectTraceForArgs(Function \&F)} iterates the function's formal
parameters, resolving each to its \texttt{DILocalVariable} via the
\texttt{DISubprogram} attached to~$F$. For each parameter whose type resolves
(through \texttt{stripDITypedefs()}) to a \texttt{DICompositeType}, the pass
calls \texttt{getStructFieldOffsets()} which:
(a)~walks the composite's member list extracting byte offsets and sizes,
(b)~computes an FNV-1a hash of the struct's name (stored at index~0 of the
offsets array for type deduplication), and
(c)~creates a module-level \texttt{ConstantArray} global.

\texttt{InjectTraceForRet(Function \&F)} uses \texttt{EscapeEnumerator} to find
all function exit points and inserts the return-value callback before each
\texttt{ReturnInst}.

Integration into the clang driver required changes to:
\texttt{CodeGenOptions.def} (two new \texttt{CODEGENOPT} flags),
\texttt{SanitizerArgs.cpp} (parsing \texttt{"trace-args"}/\texttt{"trace-ret"}
with enum values $1\ll20$ and $1\ll21$),
\texttt{BackendUtil.cpp} (wiring to \texttt{SanitizerCoverageOptions}), and
\texttt{Instrumentation.h} (adding \texttt{DataflowArgs}/\texttt{DataflowRet}
booleans).

\subsection{Kernel Backend (350~LOC)}

The kernel implementation in \texttt{kernel/kcov\_dataflow.c} (a separate file
from the existing \texttt{kcov.c}, per reviewer request) adds:
\begin{itemize}[nosep]
  \item \texttt{struct kcov\_dataflow}: owns buffer, size, sequence counter,
        and \texttt{file\_operations}.
  \item \texttt{\_\_sanitizer\_cov\_trace\_args()}: marked \texttt{notrace},
        reads struct fields via \texttt{copy\_from\_kernel\_nofault()}, writes
        the record using \texttt{READ\_ONCE}/\texttt{WRITE\_ONCE} slot reservation.
  \item \texttt{\_\_sanitizer\_cov\_trace\_ret()}: same pattern for return values.
\end{itemize}

Kconfig options \texttt{CONFIG\_KCOV\_DATAFLOW\_ARGS} and
\texttt{CONFIG\_KCOV\_DATAFLOW\_RET} (both depending on \texttt{CONFIG\_KCOV})
gate compilation.

\subsection{Build System Integration}

\texttt{scripts/Makefile.kcov} exports \texttt{CFLAGS\_KCOV\_DATAFLOW}
containing the sanitizer coverage flags. \texttt{scripts/Makefile.lib} applies
these flags when a module's Makefile declares:
\begin{lstlisting}
KCOV_DATAFLOW_my_module.o := y
\end{lstlisting}

\paragraph{Global Enablement.}
\texttt{CONFIG\_KCOV\_DATAFLOW\_INSTRUMENT\_ALL} mirrors
\texttt{CONFIG\_KCOV\_INSTRUMENT\_ALL}: when enabled, all kernel objects are
compiled with dataflow instrumentation automatically. Individual files or
directories can opt out with \texttt{KCOV\_DATAFLOW\_file.o := n} or
\texttt{KCOV\_DATAFLOW := n}. This enables whole-kernel argument/return
extraction for comprehensive fuzzing campaigns.

\paragraph{Optional Inlining Control.}
\texttt{CONFIG\_KCOV\_DATAFLOW\_NO\_INLINE} (default \texttt{n}) adds
\texttt{-fno-inline} to instrumented files for complete argument visibility.
Setting it to \texttt{y} provides full function boundary coverage at the cost
of increased stack depth and code size---suitable for targeted investigation
sessions but not continuous fuzzing.

For Rust modules, the wrapper script
\texttt{tools/kcov-dataflow/rustc-kcov-dataflow.sh} automates the three-stage
pipeline.

\subsection{User-Space Consumer}

A minimal consumer opens \texttt{/dev/kcov\_dataflow}, calls
\texttt{ioctl(KCOV\_DF\_INIT\_TRACE, buf\_size)} to allocate the buffer,
\texttt{mmap()}s it, enables recording with \texttt{ioctl(KCOV\_DF\_ENABLE)},
triggers the target kernel path, then reads records directly from the mapped
buffer before calling \texttt{ioctl(KCOV\_DF\_DISABLE)}.


\section{Evaluation}\label{sec:eval}

\subsection{Experimental Setup}

I evaluate \toolname{} on linux-next 7.1.0-rc6 (next-20260604) compiled with
my custom clang/LLVM~23 (commit \texttt{322873d}, trunk plus the dataflow
pass). Rust modules are compiled with rustc~1.98.0-nightly (2026-05-22), built
against the same LLVM~23 to share the dataflow instrumentation passes. The
kernel is configured with \texttt{CONFIG\_KASAN=y}, \texttt{CONFIG\_KCOV=y},
\texttt{CONFIG\_KCOV\_DATAFLOW\_ARGS=y}, \texttt{CONFIG\_KCOV\_DATAFLOW\_RET=y},
\texttt{CONFIG\_RUST=y}, \texttt{CONFIG\_ANDROID\_BINDER\_IPC=y}, and
\texttt{CONFIG\_SAMPLE\_RUST\_MISC\_DEVICE=y}. Tests run under virtme-ng
(QEMU/KVM) with 8~vCPUs and 1\,GB RAM.

Four selftest kernel modules (2 Rust, 2 C) exercise argument capture
correctness, and additional purpose-built modules in \texttt{findings/}
exercise distinct scenarios: FFI contract violation, silent in-bounds
corruption, 10-level deep taint propagation, binder ioctl edge cases, and
Rust core API boundary conditions (KVec, RBTree, Arc, Page, UserSlice,
credential subsystem).
Source code is available at:
kernel\footnote{\url{https://github.com/yskzalloc/linux}},
LLVM\footnote{\url{https://github.com/yskzalloc/llvm-project}}, and
rustc\footnote{\url{https://github.com/yskzalloc/rust}}.

\subsection{Correctness Verification}

Table~\ref{tab:correctness} shows that \toolname{} correctly captures
pre-corruption (ENTRY) and post-corruption (RET) argument values for all five
cases, matching \texttt{printk} ground truth.

\begin{table*}[t]
\centering
\caption{Correctness verification across five vulnerability classes.}\label{tab:correctness}
\small
\resizebox{\textwidth}{!}{%
\begin{tabular}{@{}lllll@{}}
\toprule
\textbf{Case} & \textbf{ENTRY Captured} & \textbf{RET Captured} & \textbf{drgn} & \textbf{KASAN} \\
\midrule
OOB Write (C) & \texttt{struct\{id=0x1337, buf="initial\_", size=15\}} & \texttt{struct\{id=0x1337, \_, size=32\}} & \cmark & slab-out-of-bounds \\
UAF Write (C) & \texttt{struct\{id=0, buf=poison(0xfd), size=16\}} & \texttt{struct\{0x41414141, "UAF\_CORR", 0xdead\}} & \cmark & slab-use-after-free \\
Double-Free (C) & \texttt{struct\{id=0, buf=poison(0xff), size=99\}} & \texttt{struct\{0xdf00df00, "DOUBLEFR", 0xdf\}} & \cmark & double-free \\
10-Deep Chain (C) & offset propagation: $16\to48\to24\to8$ & \texttt{commit\_write(idx=8)} OOB & \cmark & slab-out-of-bounds \\
Rust Module & \texttt{struct\{0xcafe, 0xdeadbeef.., 0xa\}}, mult=3, cookie=0x1111.. & \texttt{ret=0xefbed00023456789} & \xmark & N/A \\
\bottomrule
\end{tabular}%
}
\end{table*}

\subsection{Information Gain Over Existing Tools}

Table~\ref{tab:infogain} compares the observability provided by \toolname{}
against KASAN reports and \texttt{drgn} vmcore analysis.

\begin{table}[t]
\centering
\caption{Information gain comparison.}\label{tab:infogain}
\small
\begin{tabular}{@{}lccc@{}}
\toprule
\textbf{Observable} & \textbf{KASAN} & \textbf{drgn} & \textbf{Mine} \\
\midrule
Which argument corrupted & \xmark & \cmark~(C) & \cmark \\
Pre-corruption value & \xmark & partial & \cmark \\
Post-corruption value & \xmark & \cmark~(C) & \cmark \\
Cross-function propagation & \xmark & manual & \cmark \\
Rust module arguments & N/A & \xmark & \cmark \\
No-crash semantic bugs & \xmark & \xmark & \cmark \\
\bottomrule
\end{tabular}
\end{table}

KASAN reports the \emph{address} of a violation but not which argument carried
the bad value or which caller passed it. \texttt{drgn} can read C locals from a
vmcore but fails entirely on Rust frames compiled at \texttt{-O2} (DWARF
variable locations are elided). \toolname{} captures both C and Rust arguments
at runtime regardless of optimization level.

\subsection{Performance Overhead}

I measure overhead using a micro-benchmark that triggers 8 instrumented
functions (1--8 arguments each, 44 total callbacks per iteration) over 5{,}000
iterations in a QEMU/KVM guest, taking the best of 3 runs after warmup to
eliminate cache effects.

\begin{table}[t]
\centering
\caption{Performance overhead of \toolname{}.}\label{tab:perf}
\small
\begin{tabular}{@{}lrr@{}}
\toprule
\textbf{Configuration} & \textbf{Latency} & \textbf{Overhead} \\
\midrule
No recording (baseline) & 15.0\,$\mu$s/iter & --- \\
\toolname{} recording ON & 16.3\,$\mu$s/iter & +8.3\% \\
\bottomrule
\end{tabular}
\end{table}

The per-call cost is 1.24\,$\mu$s for 44 callbacks (8 functions $\times$
avg.\ 4.5 args + 8 returns), yielding $\sim$27\,ns per individual callback.
This cost is dominated by one \texttt{READ\_ONCE}/\texttt{WRITE\_ONCE}
atomic counter update (portable to ARM64) plus
one \texttt{copy\_from\_kernel\_nofault()} per struct field.

Critically, non-instrumented code paths incur \emph{zero} overhead---the
compiler emits no callbacks for modules without the
\texttt{KCOV\_DATAFLOW} flag. Kernel boot time is unchanged.

\paragraph{Global Enablement Overhead.}
With \texttt{CONFIG\_KCOV\_DATAFLOW\_INSTRUMENT\_ALL=y} and
\texttt{CONFIG\_KCOV\_DATAFLOW\_NO\_INLINE=n} (inlining preserved), the kernel
boots normally under virtme-ng. Table~\ref{tab:global-perf} quantifies the
overhead when all kernel objects are instrumented but recording is
\emph{not} enabled (the common idle state).

\begin{table}[htbp]
\centering
\caption{Global instrumentation overhead (recording disabled).}\label{tab:global-perf}
\small
\begin{tabular}{@{}lrrr@{}}
\toprule
\textbf{Metric} & \textbf{Baseline} & \textbf{INSTRUMENT\_ALL} & \textbf{Overhead} \\
\midrule
vmlinux .text        & 81.2\,MB & 88.9\,MB & +9.5\% \\
vmlinux .data        & 13.1\,MB & 18.9\,MB & +44\% \\
bzImage              & 33.2\,MB & 37.6\,MB & +13\% \\
Boot time            & 9.3\,s   & 15.9\,s  & +71\% \\
Syscall latency      & 20.5\,$\mu$s/iter & 47.7\,$\mu$s/iter & +133\% \\
\bottomrule
\end{tabular}
\end{table}

The .data increase (+44\%) is due to per-function \texttt{\_\_sancov\_offsets\_}
constant arrays containing struct field layouts. The syscall overhead (+133\%)
stems from per-function prologue code (alloca spills, 6-argument call setup,
branch on \texttt{kcov\_df\_enabled}) and increased icache pressure. For
comparison, basic KCOV (\texttt{trace-pc}) adds only 5--10\% with
\texttt{INSTRUMENT\_ALL} since it emits a single call per edge with no argument
setup.

Global enablement is suitable for targeted investigation sessions and
short-lived fuzzing campaigns. For continuous fuzzing (e.g., multi-day
syzkaller runs), per-module instrumentation remains recommended.

\paragraph{Comparison with Kernel Sanitizers.}
Table~\ref{tab:sanitizer-cmp} contextualizes \toolname{}'s overhead against
established kernel sanitizers. All measurements include
\texttt{CONFIG\_KASAN=y} in the baseline (the standard fuzzing configuration).

\begin{table}[htbp]
\centering
\caption{Overhead comparison with kernel sanitizers.}\label{tab:sanitizer-cmp}
\small
\begin{tabular}{@{}lrl@{}}
\toprule
\textbf{Tool} & \textbf{Syscall Overhead} & \textbf{Instruments} \\
\midrule
Basic KCOV (trace-pc)       & +5--10\%   & 1 call/edge \\
KCOV + comparisons          & +15--30\%  & edges + cmp operands \\
\toolname{} (INSTRUMENT\_ALL) & +133\%   & per-arg + per-ret callbacks \\
KASAN (generic)             & +100--200\% & every memory access \\
KCSAN                       & +100--200\% & every memory access \\
KMSAN                       & +200--400\% & every uninitialized use \\
\bottomrule
\end{tabular}
\end{table}

\noindent The overhead of \toolname{} falls within the range of a single
sanitizer (KASAN/KCSAN) and is well below KMSAN. Since fuzzing kernels
routinely stack multiple sanitizers (KASAN + KCOV + UBSAN), adding dataflow
extraction fits naturally within the accepted performance budget.
Notably, my measurements already include KASAN---the +133\% is the
\emph{additional} cost on top of an already-sanitized kernel.

\subsection{Triage Capability}

For each vulnerability case, the ENTRY/RET record provides:
(1)~the exact argument that was corrupted (identified by \texttt{arg\_idx} and
field offset), (2)~the struct field where corruption occurred (visible as a
value change between ENTRY and RET), and (3)~the call chain leading to
corruption (via sequence numbers and PC values). The 10-deep chain case
demonstrates that \toolname{} captures taint propagation across function
boundaries automatically, identifying the root cause
(\texttt{validate\_header()} missing bounds check) from the dataflow log alone.

\subsection{8-Argument Rust Kernel Module}

To verify that \toolname{} captures Rust kernel module function arguments
natively (without the post-compilation pipeline), I built \texttt{rustc}
against the custom LLVM~23 containing the trace-args/ret passes. The kernel
build system passes \texttt{-Cllvm-args=-sanitizer-coverage-trace-args}
and \texttt{-Cllvm-args=-sanitizer-coverage-trace-ret} to rustc for modules
with \texttt{KCOV\_DATAFLOW\_module.o := y}.

The \texttt{eight\_args\_rust} module defines 8 functions with 1--8 arguments
of mixed types (u64, u32, struct pointer), exercising both register-passed
(arguments 1--6 on x86-64) and stack-passed (arguments 7--8) parameters:

\begin{lstlisting}[numbers=none]
rdf_func1(0x1111111111111111)
  0x1111111111111111 = rdf_func1()
rdf_func2(0x1111111111111111, 0x22222222)
  0x33333333 = rdf_func2()
    ...
rdf_func8(0x1111.., 0x2222.., ..., struct{.x=0xaaaa, .y=0xbbbb}, 0x8888..)
  0x22222222eef05552 = rdf_func8()
\end{lstlisting}

All 8 argument values are captured individually with correct sizes. Struct
field expansion works on arg[6] (pointer to \texttt{\#[repr(C)] struct Pair}).
This confirms that Rust functions compiled with the custom rustc receive
identical instrumentation to C functions---no post-compilation pipeline
required.

\subsection{Rust FFI Contract Auditing}

The \texttt{rust\_ffi\_contract} module demonstrates \toolname{}'s primary
Rust use case: detecting contract violations at Rust-to-C FFI boundaries.
A Rust function prepares a \texttt{\#[repr(C)]} struct and passes it to a C
callee. The C function silently modifies a field that the Rust caller assumes
is immutable:

\begin{lstlisting}[numbers=none]
[ENTRY] rust_caller(struct{.id=0xcafe, .buf_ptr=0xffff8880.., .len=10})
  [ENTRY] c_callee(struct{.id=0xcafe, .buf_ptr=0xffff8880.., .len=10})
  [RET]   c_callee() = 0
[RET]   rust_caller() = 0xcafe
\end{lstlisting}

The captured records show the exact struct state at the FFI boundary.
A post-hoc verifier can compare ENTRY fields of \texttt{c\_callee} against
the Rust-side contract (e.g., \texttt{len} must not exceed buffer capacity)
and detect violations that produce no crash, no KASAN report, and no
observable symptom other than incorrect behavior downstream.

This is impossible with \texttt{ftrace} (no \texttt{-mfentry} in Rust objects),
\texttt{drgn} (DWARF elided at \texttt{-O2}), or kprobes (Rust symbol
mangling makes targeting specific functions impractical).


\section{Case Study: Runtime Boundary Auditing}\label{sec:realworld}

I applied \toolname{} to production kernel code to evaluate whether
boundary data-flow records provide actionable information beyond what
static review, KASAN, or edge coverage alone can surface.
I emphasize that the goal is \emph{demonstrating the methodology},
not claiming novel vulnerability discovery: several observations below
were subsequently clarified by Rust kernel community members as intentional behavior.

\subsection{Methodology}

For each target subsystem, I follow a four-step process:
\begin{enumerate}[nosep]
  \item Enable \texttt{KCOV\_DATAFLOW} on the target module
        (\texttt{KCOV\_DATAFLOW\_module.o := y}).
  \item Write a user-space program exercising boundary inputs (edge values,
        NULL pointers, maximum sizes).
  \item Capture the dataflow records via \texttt{mmap()} during execution.
  \item Run an \emph{automated contract verifier} that checks each captured
        $(\mathit{args}, \mathit{ret})$ pair against a contract database
        (Section~\ref{sec:contracts}).
\end{enumerate}

The key insight is that steps~3--4 require no human judgment for known
contracts: the verifier mechanically flags violations by comparing captured
values against pre/postcondition predicates. Only \emph{novel} violations
(functions without existing contracts) require manual analysis and maintainer
discussion. This shifts the workflow from ``report every observation and
wait for feedback'' to ``automatically verify known contracts; escalate only
confirmed anomalies.''

\paragraph{Why compile-time checking is insufficient.}
Rust's type system and C's \texttt{\_\_attribute\_\_} annotations enforce
structural constraints (correct types, lifetime bounds) but cannot verify
\emph{value} constraints that depend on runtime kernel state. A function
signature \texttt{alloc\_buf(size: u32)} type-checks for any \texttt{u32},
but at runtime the valid range depends on non-deterministic conditions:
available slab memory, per-task resource limits, concurrent allocations, and
RCU grace period state. These conditions are unknowable at compile time.
Boundary contract verification fills this gap: by capturing the \emph{actual}
argument values observed at runtime and checking them against value predicates
($\mathit{size} \leq \mathit{pool\_capacity}$), we verify invariants that
no static analysis or type system can express.

\emph{Critical distinction:} \toolname{} captures \emph{what the kernel
actually does} at runtime. Whether that behavior constitutes a bug requires
human judgment and maintainer confirmation---the tool provides evidence,
not verdicts.

\paragraph{Violation detection criteria.} A captured record is flagged as a
potential violation when any of the following hold:
\begin{enumerate}[nosep]
  \item An ENTRY argument value falls outside the documented or inferred
        valid range for that parameter (e.g., size $>$ buffer capacity).
  \item A RET value differs from what the precondition guarantees (e.g.,
        \texttt{from\_errno(-4096)} returns \texttt{-22} instead of
        propagating the input).
  \item A struct field changes between ENTRY and RET in a function documented
        as read-only with respect to that field.
  \item A pointer argument is \texttt{NULL} or in the \texttt{ERR\_PTR} range
        at a call site whose callee performs unconditional dereference.
\end{enumerate}
These criteria are applied offline against the captured buffer; the kernel
itself makes no policy decisions during recording.

\subsection{Observations and Kernel Community Responses}

\paragraph{Binder \texttt{SET\_MAX\_THREADS} (reported, not a bug).}
\toolname{} captured \texttt{BINDER\_SET\_MAX\_THREADS(0xffffffff)} being
accepted without bounds validation. I initially hypothesized this could
enable OOM via unbounded thread creation. Kernel maintainer Alice Ryhl
confirmed that this value does \emph{not} bypass \texttt{RLIMIT\_NPROC}:
the kernel enforces process-level thread limits independently. The recorded
dataflow correctly showed the ioctl accepting the value; the \emph{interpretation}
that this was exploitable was incorrect.

\paragraph{Binder \texttt{BC\_ENTER\_LOOPER} duplicate (reported, not a bug).}
Dataflow records showed a duplicate \texttt{BC\_ENTER\_LOOPER} command being
accepted. Rust Kernel Community review confirmed the operation is bit-level idempotent
(\texttt{looper |= BINDER\_LOOPER\_STATE\_ENTERED}) and produces no corruption.

\paragraph{\texttt{Error::from\_errno} boundary (confirmed, pre-existing TODO).}
The dataflow record captured \texttt{from\_errno(-4096)} returning \texttt{-22}
(EINVAL). The source already contains a \texttt{// TODO: Make it a WARN\_ONCE}
comment acknowledging this edge case. This demonstrates \toolname{}'s ability
to surface known-but-unresolved issues through runtime evidence rather than
source inspection.

\subsection{Rust Kernel Core API Verification}

I exercised \texttt{rust/kernel/} APIs with boundary inputs to verify
correct behavior (not to find bugs). The dataflow records confirmed:

\begin{itemize}[nosep]
  \item \textbf{KVec:} \texttt{reserve(usize::MAX)} returns
        \texttt{AllocError}; \texttt{push\_within\_capacity}
        on a full vec returns the rejected value.
  \item \textbf{RBTree:} Duplicate key insert replaces the existing value
        (documented behavior, differs from \texttt{std::HashMap::try\_insert}).
  \item \textbf{Arc:} \texttt{into\_unique\_or\_drop} with multiple refs
        correctly drops; refcount operations captured at \texttt{-O2}
        where \texttt{drgn} cannot observe them.
  \item \textbf{Page:} Cross-page boundary reads rejected
        (offset+len $>$ 4096 $\to$ EINVAL).
  \item \textbf{UserSlice:} NULL, kernel-range, and unmapped addresses all
        return EFAULT correctly.
\end{itemize}

\noindent These results demonstrate that \toolname{} is equally valuable for
\emph{building confidence in correct implementations} as for identifying
potential issues. The dataflow records provide reproducible evidence that
APIs handle boundary inputs as intended---evidence that static review alone
cannot provide at \texttt{-O2}.

\subsection{Correction Note}

In version~1 of this paper, I identified several potential vulnerabilities
in the Android Binder driver and Rust kernel core APIs, assigning severity
labels (High, Critical) prior to maintainer verification. Following review
and discussion with the Linux kernel maintainers, these specific findings
were determined to be not exploitable under standard configurations or
addressed by existing kernel mitigations:

\begin{itemize}[nosep]
  \item \texttt{BINDER\_SET\_MAX\_THREADS(0xffffffff)}: not exploitable---
        \texttt{RLIMIT\_NPROC} enforcement is independent and prevents
        unbounded thread creation regardless of the stored value.
  \item \texttt{BC\_ENTER\_LOOPER} duplicate: not a corruption vector---
        the operation is bit-level idempotent.
  \item Rust API behavioral observations (e.g., \texttt{RBTree} duplicate
        key semantics): documented intentional behavior, not bugs.
\end{itemize}

\noindent Consequently, these findings have been removed from this version.
The system design, methodology, and the runtime boundary auditing framework
presented in the preceding sections remain valid. Verified findings will be
incorporated step by step in subsequent versions as maintainer review
completes.

\section{Why ftrace Cannot Trace Rust Functions}\label{sec:ftrace}

A natural question is whether existing tracing infrastructure (\texttt{ftrace},
\texttt{kprobes}) could achieve similar results. The answer is no, for
fundamental architectural reasons.

\paragraph{ftrace requires \texttt{-mfentry}.} The C compiler emits a
\texttt{call \_\_fentry\_\_} prologue when passed \texttt{-pg -mfentry}. The
kernel's build system passes this flag to \texttt{clang}/\texttt{gcc} for C
files. However, \texttt{rustc} does not understand \texttt{-mfentry}---it has
its own codegen pipeline. No \texttt{\_\_fentry\_\_} call is emitted in Rust
object files, so ftrace has nothing to patch at runtime.

\paragraph{ftrace captures no arguments.} Even for C functions where ftrace
works, it only reports ``function X was entered/exited.'' Obtaining argument
values requires writing a custom kprobe handler or eBPF program that reads
specific registers---which requires knowing the exact calling convention and
register assignment for each function.

\paragraph{Rust symbol mangling.} Rust function names are mangled into
unreadable strings (e.g., \texttt{\_RNvMs3\_NtCsgHKxaxYGgF6\_11rust\_binder...}).
Setting ftrace filters or kprobe targets requires these exact names.

\paragraph{\toolname{} operates at LLVM IR level.} My pass runs on LLVM IR
\emph{before} codegen---the common representation shared by both \texttt{clang}
(C) and \texttt{rustc} (Rust). At this level, function arguments are explicit
SSA values with debug metadata, struct types have \texttt{DICompositeType}
annotations, and the optimizer has not yet destroyed variable locations. The
inserted callbacks survive optimization because they are function calls with
side effects.

\begin{table}[t]
\centering
\caption{Comparison: ftrace vs.\ \toolname{} for Rust kernel code.}\label{tab:ftrace}
\small
\begin{tabular}{@{}lcc@{}}
\toprule
\textbf{Capability} & \textbf{ftrace} & \textbf{\toolname{}} \\
\midrule
Trace Rust functions & \xmark & \cmark \\
Capture argument values & \xmark & \cmark \\
Struct field decomposition & \xmark & \cmark \\
Zero source modification & \cmark & \cmark \\
Works at \texttt{-O2} & \cmark & \cmark \\
Per-task isolation & \xmark & \cmark \\
\bottomrule
\end{tabular}
\end{table}

\section{Boundary Contract Runtime Verification}\label{sec:contracts}

I observe that \toolname{} enables a new form of \emph{runtime verification}
for kernel functions: verifying pre/postcondition contracts at function
boundaries without source modification, for any function the monitored process
executes---including functions the developer did not anticipate.

\subsection{Functions as Black-Box Subsystems}

Every kernel function call is a boundary between a \emph{caller} (who provides
inputs) and a \emph{callee} (a black box from the caller's perspective):
$$\text{Caller} \xrightarrow{[\mathit{arg}_1, \ldots, \mathit{arg}_n]}
  f_{\text{callee}} \xrightarrow{[\mathit{ret}]} \text{Caller}$$

The caller cannot observe what happens inside the callee. It can only choose
arguments (precondition) and observe the return value (postcondition).
\toolname{} captures \emph{both sides} of this boundary, making the black box
transparent without modifying it.

\subsection{Formal Boundary Contracts}

Define a \emph{Boundary Contract} $\mathcal{C}(f) = (\mathcal{P},
\mathcal{Q})$ for function $f$:
\begin{itemize}[nosep]
  \item $\mathcal{P}(\mathit{args})$: precondition predicate over input arguments
  \item $\mathcal{Q}(\mathit{args}, \mathit{ret})$: postcondition predicate
        over return value given inputs
\end{itemize}

\noindent\textbf{Example:} For \texttt{binder\_alloc\_buf(alloc, data\_size,
offsets\_size, is\_async)}:
\begin{align*}
\mathcal{P}&: \mathit{data\_size} \leq 2^{20} \wedge
              \mathit{offsets\_size} \leq 2^{16} \\
\mathcal{Q}&: \mathit{ret} = 0 \implies
              \mathit{alloc.buffer} \neq \texttt{NULL}
\end{align*}

A \emph{contract violation} occurs when the caller passes arguments violating
$\mathcal{P}$ (caller bug) or the callee returns values violating $\mathcal{Q}$
given valid inputs (callee bug).

\subsection{The Unknown Kernel Path Problem}

When a user-space process calls \texttt{write(binder\_fd, data, len)}, it
triggers a chain of kernel functions:
\begin{lstlisting}[numbers=none]
sys_write -> vfs_write -> binder_write ->
  binder_thread_write -> binder_transaction ->
  binder_alloc_buf -> ...
\end{lstlisting}

The user cannot predict which functions will execute. Linux
RV~\cite{dealmeida2019rv} requires
pre-defining which events to monitor. \toolname{} captures \emph{every}
function boundary automatically---the offline verifier then checks each
captured $(\mathit{args}, \mathit{ret})$ pair against a contract database.

\subsection{Comparison with Linux RV}

\begin{table}[t]
\centering
\caption{Linux RV vs.\ \toolname{} contract verification.}\label{tab:rv}
\small
\begin{tabular}{@{}lcc@{}}
\toprule
\textbf{Dimension} & \textbf{Linux RV} & \textbf{\toolname{}} \\
\midrule
Scope & Pre-defined events & Any function executed \\
Model & Automata on sequences & Predicates on values \\
Instrumentation & Explicit tracepoints & Zero source modification \\
Catches & Ordering violations & Value violations \\
Non-deterministic state & No (static model) & Yes (runtime values) \\
Unknown functions & Cannot verify & Captures automatically \\
Rust support & No & Yes \\
Specification & Hardcoded in kernel & External database \\
Compile-time sufficiency & N/A & No (value depends on runtime) \\
Overhead & Near zero (event-driven) & $<$10\% per-module \\
\bottomrule
\end{tabular}
\end{table}

Table~\ref{tab:rv} summarizes the differences. Linux RV verifies
\emph{anticipated} event orderings (e.g., ``sleep must not occur in atomic
context''). \toolname{} verifies \emph{arbitrary} value contracts (e.g.,
``\texttt{data\_size} must be $\leq 2^{20}$'') on functions the developer
never explicitly instrumented.

\subsection{Concrete Example: Detecting the errno Bug}

Contract for \texttt{Error::from\_errno(errno)}:
$$\mathcal{P}: -4095 \leq \mathit{errno} \leq -1$$
$$\mathcal{Q}: \mathit{ret.code} = \mathit{errno}$$

\toolname{} capture:
\begin{lstlisting}[numbers=none]
[ENTRY] pc=from_errno arg[0] = -4096  // P violated
[RET]   pc=from_errno ret = -22       // Q violated
\end{lstlisting}

Verifier output: precondition violated ($-4096 \notin [-4095, -1]$) and
postcondition violated ($\mathit{ret} = -22 \neq -4096$). Linux RV cannot
catch this: no tracepoint exists inside \texttt{from\_errno}, and it is a Rust
function invisible to C-based tracing infrastructure.

\subsection{Struct Member Contract Verification}\label{sec:struct-contract}

Beyond scalar argument bounds, \toolname{}'s automatic struct field expansion
enables \emph{struct member contracts}: predicates over relationships between
fields that must hold for a function's internal logic to operate correctly.

\paragraph{The problem.} Most kernel functions receive struct pointers whose
fields encode complex state (flags, sizes, pointers, reference counts). The
function's internal logic assumes these fields satisfy certain relationships
(e.g., \texttt{buf\_len $\geq$ data\_offset + data\_size}). These
relationships are typically enforced by conditional checks inside the function,
but:
\begin{itemize}[nosep]
  \item Checks may be missing for uncommon field combinations.
  \item Concurrent modification may invalidate a field between the check and
        its use (TOCTOU).
  \item The caller may construct a struct in an inconsistent state that passes
        individual field validation but violates cross-field invariants.
\end{itemize}

Compile-time type systems cannot detect these: the struct is well-typed
regardless of whether its fields are internally consistent. The correct state
depends on runtime conditions (which allocator path was taken, whether a lock
is held, what RCU grace period is active).

\paragraph{Formalization.} Extend the boundary contract to include a
\emph{struct consistency predicate} $\mathcal{S}$ over fields of a struct
parameter:
$$\mathcal{S}(\mathit{fields}): \bigwedge_i \phi_i(f_1, f_2, \ldots, f_n)$$

\noindent\textbf{Example:} For \texttt{smb2\_create\_req *req}:
\begin{align*}
\mathcal{S}&: \mathit{req.NameLength} \leq \mathit{req.BufferLength} \\
           &\wedge\ \mathit{req.CreateContextsOffset} + \mathit{req.CreateContextsLength}
              \leq \mathit{req.StructureSize}
\end{align*}

\paragraph{Detection mechanism.} At function entry, \toolname{} captures all
struct fields simultaneously (one atomic record per parameter). The offline
verifier evaluates $\mathcal{S}$ over the captured field values. A violation
means the function received a struct whose internal state is inconsistent with
what its logic requires---either the caller constructed it incorrectly, or
concurrent modification corrupted it between construction and use.

\paragraph{Concrete example.} Consider a kernel network handler:
\begin{lstlisting}[numbers=none]
[ENTRY] process_msg(struct{.hdr_len=20, .payload_len=4096,
                           .total_len=100, .flags=0x3})
\end{lstlisting}
Contract: $\mathit{hdr\_len} + \mathit{payload\_len} \leq \mathit{total\_len}$.
Captured: $20 + 4096 = 4116 > 100$ --- violation. The function will either
read out-of-bounds or silently produce incorrect output. No crash occurs
(the buffer may be large enough by coincidence), no KASAN report fires,
and edge coverage sees the same path regardless of field values.

\paragraph{Why this matters.} These bugs are among the hardest to find:
\begin{itemize}[nosep]
  \item They produce no immediate crash (data is read, just incorrectly).
  \item Sanitizers do not fire (the access is technically in-bounds of the
        allocation, just semantically wrong).
  \item Static analysis cannot reason about runtime struct state.
  \item They often manifest only under specific timing or workload conditions.
\end{itemize}
Struct member contract verification with \toolname{} is the only automated
method that can detect these violations at runtime with zero source annotation,
by comparing the actual field values against the cross-field predicates that
the function's conditional logic implicitly assumes.

\begin{figure}[t]
\centering
\begin{tikzpicture}[
  node distance=0.4cm and 0.6cm,
  box/.style={rectangle, draw, rounded corners, minimum width=2.4cm, minimum height=0.6cm, font=\scriptsize, align=center},
  struct/.style={rectangle, draw, thick, fill=yellow!10, minimum width=2.8cm, font=\scriptsize, align=left, inner sep=4pt},
  check/.style={diamond, draw, fill=orange!10, minimum width=1.2cm, minimum height=0.8cm, font=\scriptsize, align=center, inner sep=1pt},
  arrow/.style={-{Stealth[length=2mm]}, thick},
  note/.style={font=\tiny\itshape, text=gray}
]
% Function entry
\node[box, fill=blue!8] (entry) {$f(\texttt{struct}\ S\ *s, \mathit{globals})$};

% Struct state
\node[struct, below left=0.8cm and 0.3cm of entry] (sinput) {%
  \texttt{s->len} = 4096\\
  \texttt{s->cap} = 256\\
  \texttt{s->flags} = 0x3};
\node[note, above=0.05cm of sinput] {Input struct state};

% Global state
\node[struct, below right=0.8cm and 0.3cm of entry] (gstate) {%
  \texttt{pool\_free} = 128\\
  \texttt{rcu\_epoch} = 7\\
  \texttt{lock\_held} = 0};
\node[note, above=0.05cm of gstate] {Runtime global state};

% Contract check
\node[check, below=2.2cm of entry] (verify) {$\mathcal{S} \wedge \mathcal{G}$\\valid?};

% Outcomes
\node[box, fill=green!15, below left=0.8cm and 0.2cm of verify] (pass) {$\checkmark$ Consistent\\safe to proceed};
\node[box, fill=red!15, below right=0.8cm and 0.2cm of verify] (fail) {$\times$ Violation\\flag for review};

% Arrows
\draw[arrow] (entry.south) -- ++(0,-0.3) -| (sinput.north);
\draw[arrow] (entry.south) -- ++(0,-0.3) -| (gstate.north);
\draw[arrow] (sinput.south) |- (verify.west);
\draw[arrow] (gstate.south) |- (verify.east);
\draw[arrow] (verify.south) -- ++(0,-0.3) -| (pass.north) node[pos=0.3, left, note] {yes};
\draw[arrow] (verify.south) -- ++(0,-0.3) -| (fail.north) node[pos=0.3, right, note] {no};

% Capture annotation
\draw[decorate, decoration={brace, amplitude=3pt, mirror, raise=15pt}, gray]
  (sinput.south west) -- (gstate.south east)
  node[midway, below=18pt, font=\tiny\itshape, text=black] {captured by \toolname{} at function entry};
\end{tikzpicture}
\caption{Struct member and global state verification at function entry.
\toolname{} captures the struct fields and observable global state;
the offline verifier evaluates the consistency predicate
$\mathcal{S} \wedge \mathcal{G}$.}\label{fig:struct-verify}
\end{figure}

\paragraph{Formal model.}
Let a function $f$ receive a struct parameter $s$ with fields
$(f_1, \ldots, f_n)$ and access global kernel state $G = (g_1, \ldots, g_m)$.
Define the \emph{functional verification predicate}:
$$\mathcal{V}(f) = \mathcal{S}(f_1, \ldots, f_n) \wedge
  \mathcal{G}(f_1, \ldots, f_n, g_1, \ldots, g_m)$$
where:
\begin{itemize}[nosep]
  \item $\mathcal{S}$: \emph{structural consistency} --- cross-field invariants
        that must hold regardless of runtime state
        (e.g., $\mathit{len} \leq \mathit{cap}$,
        $\mathit{offset} + \mathit{size} \leq \mathit{total}$).
  \item $\mathcal{G}$: \emph{state-dependent validity} --- predicates that
        relate struct fields to runtime kernel state
        (e.g., $\mathit{alloc\_size} \leq \mathit{pool\_free}$,
        $\mathit{flags} \wedge \mathtt{LOCKED} \implies \mathit{lock\_held} = 1$).
\end{itemize}

\noindent A \emph{violation} $\neg\mathcal{V}(f)$ at function entry means the
function's internal conditional logic will either:
\begin{enumerate}[nosep]
  \item Take an unanticipated branch (missing check $\to$ silent corruption), or
  \item Produce semantically incorrect output (wrong branch $\to$ logic bug).
\end{enumerate}

\paragraph{What \toolname{} captures.} For struct fields, all values are
captured via \texttt{DICompositeType} expansion at entry (Section~\ref{sec:llvm-pass}).
For globals accessed by the function, they appear as arguments if passed
explicitly, or can be captured at a wrapper boundary. The key insight is that
\toolname{} provides a \emph{snapshot} of the state at the exact moment the
function is invoked---before any internal check can mask the inconsistency.

\paragraph{Classification of violations.}
\begin{center}
\small
\begin{tabular}{@{}lll@{}}
\toprule
\textbf{Violation type} & \textbf{Predicate} & \textbf{Consequence} \\
\midrule
Cross-field overflow & $\mathit{off} + \mathit{len} > \mathit{cap}$ & OOB read/write \\
Flag-state mismatch & $\mathit{flags} \wedge F \neq 0 \wedge \neg\mathit{ready}$ & Use-before-init \\
Size-resource conflict & $\mathit{req\_sz} > \mathit{avail}$ & Allocation failure \\
Epoch staleness & $\mathit{s.epoch} < \mathit{global\_epoch}$ & Stale data use \\
\bottomrule
\end{tabular}
\end{center}

\subsection{Toward Post-Mortem Formal Verification}\label{sec:postmortem}

Classical formal verification proves properties for \emph{all} inputs
statically, but is intractable for 30M~LOC kernels. I observe that
\toolname{}'s structured output enables a complementary approach:
\emph{post-mortem contract verification}---mechanically checking recorded
execution against formal specifications after the fact.

\paragraph{Key Insight.} The ring buffer produces a complete sequence of
$\langle \mathit{PC}, \mathit{args}, \mathit{ret} \rangle$ tuples for every
instrumented function the process executed. This is a \emph{finite execution
trace} in the formal methods sense---the input to a runtime verification
monitor.

\paragraph{What This Enables.}
\begin{enumerate}[nosep]
  \item \textbf{Offline model checking:} Replay the buffer against temporal
        logic specifications (e.g., LTL: ``every \texttt{alloc} is eventually
        followed by \texttt{free} with the same pointer'').
  \item \textbf{Contract refinement:} Automatically infer likely contracts from
        many executions (Daikon-style invariant detection on kernel functions),
        then flag violations in new runs.
  \item \textbf{Differential verification:} Compare argument/return value
        distributions across kernel versions to detect semantic regressions
        that pass all tests but violate implicit contracts.
\end{enumerate}

\paragraph{Soundness Boundary.} This is \emph{not} full formal verification:
it only covers paths actually exercised. However, combined with a fuzzer that
maximizes path diversity, it approaches bounded model checking---verifying
contracts over all \emph{reachable} states discovered during exploration.
The structured, machine-readable output (unlike \texttt{printk} logs) makes
this amenable to automated reasoning tools without manual log parsing.

\section{Discussion \& Limitations}\label{sec:discussion}

\subsection{Runtime-Originated Kernel Code Auditing}\label{sec:workflow}

I define \emph{runtime-originated auditing} as the methodology in which the
analysis target is not the source code itself, but the dataflow records captured
during actual kernel execution triggered by a user-space operation. This inverts
the conventional workflow and eliminates the speculative gap inherent in
source-only review.

\begin{table}[t]
\centering
\caption{Source-speculative vs.\ evidence-based auditing.}\label{tab:methodology}
\small
\begin{tabular}{@{}lll@{}}
\toprule
\textbf{Dimension} & \textbf{Without \toolname{}} & \textbf{With \toolname{}} \\
\midrule
Input & Source code & Captured records \\
Evidence & ``Path appears reachable'' & PC + arg values at runtime \\
Validation & Crash or sanitizer required & Record is the evidence \\
Reproducibility & Analyst-dependent & Deterministic artifact \\
Scope & Anticipated paths only & Every boundary executed \\
\bottomrule
\end{tabular}
\end{table}

\begin{figure*}[t]
\centering
\begin{tikzpicture}[
  node distance=0.35cm,
  box/.style={rectangle, draw, rounded corners, minimum width=2.8cm, minimum height=0.65cm, font=\scriptsize, align=center},
  arrow/.style={-{Stealth[length=2mm]}, thick},
  title/.style={font=\scriptsize\bfseries}
]
% Left: Existing (KASAN + KCOV + Fuzzer)
\node[title, text=red!70!black] (lt) at (0,0) {Existing: KASAN + KCOV + Fuzzer};
\node[box, fill=red!5, below=0.3cm of lt] (l1) {Fuzzer mutates syscalls};
\node[box, fill=red!5, below=of l1] (l2) {KCOV: edge coverage\\(which paths hit?)};
\node[box, fill=red!5, below=of l2] (l3) {KASAN: memory error\\(crash or no signal)};
\node[box, fill=red!5, below=of l3] (l4) {Coverage saturates\\on stateful subsystems};
\draw[arrow] (l1) -- (l2);
\draw[arrow] (l2) -- (l3);
\draw[arrow] (l3) -- (l4);
\node[font=\tiny\itshape, text=gray, below=0.1cm of l4] {blind to argument values};

% Middle: kcov-dataflow
\node[title, text=green!50!black] (mt) at (5.5,0) {\toolname{}};
\node[box, fill=green!5, below=0.3cm of mt] (m1) {Same syscall triggers};
\node[box, fill=green!8, below=of m1] (m2) {Capture (PC, args, ret)\\at every boundary};
\node[box, fill=green!8, below=of m2] (m3) {Struct fields auto-expanded\\via DICompositeType};
\node[box, fill=green!5, below=of m3] (m4) {Value-aware feedback\\breaks saturation};
\draw[arrow] (m1) -- (m2);
\draw[arrow] (m2) -- (m3);
\draw[arrow] (m3) -- (m4);
\node[font=\tiny\itshape, text=gray, below=0.1cm of m4] {observes what drove the path};

% Right: Combined
\node[title, text=blue!70!black] (rt) at (11,0) {Combined Workflow};
\node[box, fill=blue!5, below=0.3cm of rt] (r1) {Fuzzer + edge coverage\\(explore structure)};
\node[box, fill=blue!5, below=of r1] (r2) {\toolname{}: argument values\\(explore state space)};
\node[box, fill=blue!5, below=of r2] (r3) {KASAN/KMSAN: detect\\corruption when triggered};
\node[box, fill=blue!8, below=of r3] (r4) {Root-cause from dataflow:\\which arg $\to$ which field};
\draw[arrow] (r1) -- (r2);
\draw[arrow] (r2) -- (r3);
\draw[arrow] (r3) -- (r4);
\node[font=\tiny\itshape, text=gray, below=0.1cm of r4] {full observability stack};

% Dividers
\draw[dashed, gray] (2.8,0.2) -- (2.8,-4.5);
\draw[dashed, gray] (8.2,0.2) -- (8.2,-4.5);
\end{tikzpicture}
\caption{Left: existing kernel fuzzing stack (edge-blind to values).
Middle: \toolname{} adds boundary data-flow extraction.
Right: combined workflow provides full observability from mutation guidance
through root-cause analysis.}\label{fig:workflow}
\end{figure*}

\begin{figure}[t]
\centering
\begin{tikzpicture}[
  node distance=0.4cm,
  box/.style={rectangle, draw, rounded corners, minimum width=3cm, minimum height=0.7cm, font=\scriptsize, align=center},
  arrow/.style={-{Stealth[length=2mm]}, thick},
]
% Top: capture
\node[box, fill=yellow!15] (cap) {User action $\to$ \toolname{} capture\\(PC, args, ret, struct fields)};

% Split
\node[box, fill=green!10, below left=0.8cm and -0.5cm of cap] (audit) {Security Audit\\identify contract violations};
\node[box, fill=red!8, below right=0.8cm and -0.5cm of cap] (fuzz) {Guided Fuzzing\\constrain mutation ranges};

% Bottom
\node[box, fill=blue!8, below=1.5cm of cap, yshift=-1cm] (result) {Confirmed finding with\\reproducible evidence};

\draw[arrow] (cap.south) -- ++(0,-0.3) -| (audit.north);
\draw[arrow] (cap.south) -- ++(0,-0.3) -| (fuzz.north);
\draw[arrow] (audit.south) -- ++(0,-0.3) -| (result.north);
\draw[arrow] (fuzz.south) -- ++(0,-0.3) -| (result.north);
\end{tikzpicture}
\caption{Dual utility: the same captured records serve both manual security
auditing and automated fuzzing guidance, converging on confirmed findings
with reproducible evidence.}\label{fig:dual}
\end{figure}

\paragraph{Key property.} The dataflow records are a \emph{deterministic
artifact}: given the same kernel build, the same user action produces the same
records. Any party can independently reproduce and verify the analysis,
satisfying the reproducibility requirement for empirical security research.

\paragraph{Dual utility.} The same records serve both (1)~security auditing
(identifying argument values that violate API contracts) and (2)~guided fuzzing
(constraining mutation to value ranges observed at deep kernel code paths,
eliminating blind 64-bit enumeration).

\paragraph{Implication for AI-assisted auditing.} Without runtime evidence,
AI-based code review produces hypotheses that may be incorrect (as demonstrated
by my initial binder claims refuted by maintainers). With \toolname{}, the AI
analyzes observed records rather than speculating about behavior---findings are
grounded in physical execution regardless of reasoning quality.

\paragraph{Debug Info Requirement.} Full struct field expansion requires
\texttt{CONFIG\_DEBUG\_INFO} (\texttt{-g}); without it, only scalar argument
values are captured. This is not a practical limitation for kernel fuzzing
(KASAN already requires debug info).

\paragraph{Rust Pipeline.} Dataflow tracking of Rust kernel modules requires
building \texttt{rustc} against the custom LLVM~23 that contains the
trace-args/ret passes. The post-compilation pipeline (Path~A) serves as a
proof-of-concept demonstrating language-agnosticism but cannot instrument
modules built through Kbuild. I propose adding native
\texttt{-Zsanitizer-coverage=trace-args} support to upstream rustc as future
work to eliminate the custom build requirement.

\paragraph{Buffer Overflow.} If instrumented code generates more events than
the buffer capacity, events are silently dropped. This is by design---blocking
would introduce deadlocks in interrupt context.

\paragraph{Interrupt and NMI Context.} Instrumentation callbacks are
restricted to process context via an \texttt{in\_task()} check: any call
originating from a softirq, hardirq, or NMI handler is silently discarded.
This is stronger than the per-context filtering used by mainline KCOV
(which allows softirq tracing via \texttt{kcov\_remote\_start}); I chose
the stricter policy because the dataflow callback performs pointer
dereferences (\texttt{copy\_from\_kernel\_nofault}) whose latency is
unacceptable in hard-interrupt paths, and because per-task buffer
interleaving from nested contexts would corrupt record boundaries.
Bugs triggered exclusively from interrupt context are therefore invisible
to \toolname{}; however, such bugs are also unreachable by KCOV-guided
fuzzers, which exercise kernel code through syscalls in process context.

\paragraph{Fuzzer Integration (Future Work).} Integrating \toolname{} with
syzkaller requires modifying the executor to parse the \texttt{mmap}'d
\texttt{kcov\_dataflow} buffer and hash argument values into the coverage
bitmap. I envision a composite feedback signal: unique
$(\mathit{PC}, \mathit{arg\_hash})$ pairs provide finer-grained state
discrimination than unique edges alone.

\paragraph{Upstream Path.} The LLVM pass has been submitted as a pull
request\footnote{\url{https://github.com/llvm/llvm-project/pull/201410}} with
an RFC on LLVM
Discourse\footnote{\url{https://discourse.llvm.org/t/rfc-sanitizercoverage-add-fsanitize-coverage-trace-args-trace-ret/91026}};
LLVM maintainers have requested a formal RFC per the LLVM RFC
process\footnote{\url{https://llvm.org/docs/RFCProcess.html}} given the
complexity of extending \sancov with a new instrumentation mode.
The kernel patch series is under review on
LKML\footnote{\url{https://lore.kernel.org/all/20260603-kcov-dataflow-next-20260603-v2-0-fee0939de2c4@est.tech/}}.

\section{Related Work}\label{sec:related}

\paragraph{Coverage-Guided Kernel Fuzzing.}
syzkaller~\cite{vyukov2015syzkaller} and kAFL~\cite{schumilo2017kafl} use
edge/block coverage as feedback. DIFUZE~\cite{corina2017difuze} adds
interface-aware generation but no runtime data-flow.
MORPHUZZ~\cite{bulekov2022morphuzz} fuzzes virtual devices without data-flow
context. Healer~\cite{sun2021healer} learns syscall dependencies from coverage
to improve sequence construction, yet remains blind to argument values.
SyzVegas~\cite{wang2021syzvegas} applies multi-armed bandit algorithms to
syzkaller's mutation strategy but still uses edge coverage as the sole reward
signal. ACTOR~\cite{fleischer2023actor} targets concurrency bugs via
schedule-aware fuzzing; its feedback is scheduling order, not data-flow.
Moonshine~\cite{pailoor2018moonshine} distills seeds from system call traces
to improve initial coverage but cannot guide mutation once coverage saturates.
HEALER~\cite{sun2021healer} and Rtkaller~\cite{ma2022rtkaller} learn
relation models between syscalls to construct deeper call sequences, yet both
plateau once the structural path space is exhausted.
Google's AIxCC kernel fuzzing approach~\cite{google2024aixcc} combines AFL
with KCOV and KASAN in a QEMU harness, reporting that ``the vulnerability
cannot be easily found even after hours of fuzzing, unless provided with seed
inputs close to the solution''---confirming that edge coverage alone saturates
on complex kernel targets.
All of the above rely solely on control-flow coverage; none capture the
argument values that drive state-dependent behavior.

\paragraph{Data-Flow Guided Fuzzing.}
Payer~\cite{payer2019dataflow} argues for data-flow-guided fuzzing but provides
no kernel implementation. Angora~\cite{chen2018angora} uses byte-level taint for
user-space programs at 10$\times$ overhead.
GREYONE~\cite{gan2020greyone} combines taint with coverage but targets
user-space only. None achieve the $<$5\% overhead required for continuous kernel
fuzzing.

\paragraph{Kernel Instrumentation.}
KCOV~\cite{vyukov2016kcov} provides edge coverage---my foundation.
\texttt{ftrace}/\texttt{kprobes} offer dynamic per-function tracing at high
overhead. eBPF is programmable but restricted in context and cannot perform
struct field expansion from debug metadata.

\paragraph{Post-Mortem Analysis.}
\texttt{drgn}~\cite{drgn} enables programmable vmcore analysis but requires a
crash and fails on optimized Rust code. \texttt{kdump}/\texttt{crash} provide
full memory dumps but require manual analysis and miss silent bugs.

\medskip\noindent
\toolname{} is the first system combining compile-time data-flow instrumentation
with a production-grade kernel transport, supporting both C and Rust with
$<$10\% per-module overhead.

\section{Conclusion}\label{sec:conclusion}

I presented \toolname{}, bridging the semantic gap between coverage-guided
fuzzing and data-flow analysis for OS kernels. My key insight is that function
boundaries are natural observation points where data-flow context is both cheap
to capture and highly informative for both automated exploration and human
triage.

The system requires no source annotation, operates safely in process context
with explicit rejection of interrupt and NMI callbacks,
and---uniquely---supports Rust kernel modules via both a post-compilation
pipeline and native instrumentation through a custom-built rustc, providing
the only known runtime methods for capturing Rust function arguments
given \texttt{-O2} DWARF elision.

I demonstrated correctness on five vulnerability classes across C and Rust
kernel modules, showing that boundary data-flow context provides information
unavailable through any combination of KASAN, drgn, and edge coverage.
Source code is available at \anonrepo{}.

%Bibliography
\bibliographystyle{unsrt}
\bibliography{../shared/references}

\appendix
\section{LLVM Pass: InjectTraceForArgs Pseudocode}\label{app:pass}

\begin{algorithmic}[1]
\Function{InjectTraceForArgs}{$F$}
  \State $\mathit{DI} \gets F.\mathrm{getSubprogram()}$
  \If{$\mathit{DI} = \mathrm{null}$} \Return \EndIf
  \For{each parameter $P_i$ of $F$}
    \State $T \gets \mathrm{stripTypedefs}(\mathit{DI}.\mathrm{getType}(i))$
    \If{$T$ is \texttt{DICompositeType}}
      \State $\mathit{offsets} \gets \mathrm{getStructFieldOffsets}(T)$
      \State $\mathit{hash} \gets \mathrm{FNV1a}(T.\mathrm{getName()})$
      \State Create global: $[\mathit{hash}, \mathit{off}_0, \mathit{sz}_0, \ldots]$
    \Else
      \State $\mathit{offsets} \gets \mathrm{null}$; $\mathit{nfields} \gets 0$
    \EndIf
    \State $\mathit{ptr} \gets$ spill $P_i$ to alloca if scalar, else use directly
    \State Insert call: \texttt{trace\_args(pc, $i$, sizeof($P_i$), ptr, offsets, nfields)}
  \EndFor
\EndFunction
\end{algorithmic}

\section{Ring Buffer Record Format}\label{app:ringbuf}

Each record occupies $3 + N$ quadwords (64-bit each), where $N$ is the number of fields (or 1 for scalar values when $N=0$):

\begin{center}
\small
\begin{tabular}{|c|l|p{10.7cm}|}
\hline
\textbf{Offset} & \textbf{Field} & \textbf{Description} \\
\hline
0 & \texttt{type\_and\_seq} & Type marker (bits [31:28]: \texttt{0xE}=entry, \texttt{0xF}=return) | seq (bits [23:0]) \\
1 & \texttt{pc} & Instrumented function address \\
2 & \texttt{meta} & arg\_idx [63:56] | size [55:48] | pointer [47:0] \\
3..3+N & \texttt{field\_val[0..N]} & Field or scalar values via \texttt{copy\_from\_kernel\_nofault} \\
\hline
\end{tabular}
\end{center}

\section{Rust Post-Compilation Pipeline}\label{app:rust}

\begin{lstlisting}[language=bash,caption={Rust module instrumentation pipeline}]
# 1. Emit LLVM IR with debug info
rustc --edition 2021 --emit=llvm-ir -g \
  -o module.ll module.rs --crate-type lib

# 2. Instrument with my opt (uses custom LLVM 23)
opt -passes=sancov-module \
  -sanitizer-coverage-level=3 \
  -sanitizer-coverage-trace-args \
  -sanitizer-coverage-trace-ret \
  -S module.ll -o module_inst.ll

# 3. Compile to object
llc -filetype=obj -relocation-model=static \
  -code-model=kernel module_inst.ll -o module.o

# 4. Link as normal kernel module
\end{lstlisting}

\section{Ioctl Interface Specification}\label{app:ioctl}

\begin{lstlisting}[language=C,caption={Ioctl definitions}]
#define KCOV_DF_INIT_TRACK  _IOR('d', 1, unsigned long)
  /* = 0x80086401 */
#define KCOV_DF_ENABLE      _IO('d', 100)
  /* = 0x6464 */
#define KCOV_DF_DISABLE     _IO('d', 101)
  /* = 0x6465 */
\end{lstlisting}

\end{document}
